
\needspace{4\baselineskip}
\color{uniblue}{\subsection{Предупредительная сигнализация устройства (ПС)\label{ps:punkt}}}
\color{black}

\begin{enumerate}[label=\arabic{section}.\arabic{subsection}.\arabic*, labelsep=4pt, leftmargin=0pt, itemindent=57pt]

\item	
Предупредительная сигнализация предназначена для информирования оперативного персонала при отклонениях от нормального режима работы оборудования энергообъекта.

\item
В устройстве предупредительная сигнализация представлена функцией «ПС».

\item
Функция принимает сигналы срабатывания защитных и сервисных функций устройства и формирует следующие сигналы:
\begin{itemize}
	\item импульсный сигнал предупредительной сигнализации «ПС: Пуск (имп)»;
	\item длительный сигнал предупредительной сигнализации «ПС: Пуск».
\end{itemize}

\item
Введенный ключ «Вывод терминала» блокирует формирование выходных сигналов.

\item
Сигналы из логики СС:«Вых.цепи разобраны», «БИ выведены», «ОТ сигн», «Неисправность КА», «Прев. вр.пер. КА», «Общ.внеш.сигнал» контролируются соответствующими уставками SGF1\ldots SGF6.

\item
Инвертированное состояние дискретных сигналов при введенной уставке SGF7 «Полож.пер.ав.дебл. 1», SGF8 «Полож.пер.ав.дебл. 2»  формирует работу логики ПС.

\item
Выходные сигналы функционального блока могут быть использованы для дальнейшего конфигурирования (например, на выходные реле) и применяться в цепях участковых и индивидуальных табло, сигнальных ламп, сигнальных реле, звуковой сигнализации и пр.

\item
Логика работы функции выполнена в соответствии с рисунком приложения \ref{app:ps}. В таблице \ref{ps:tbl1} приведены параметры, необходимые для настройки функции. 

\end{enumerate}